\documentclass[11pt]{article}

\usepackage[T1]{fontenc}
\usepackage[margin=1in]{geometry}
\usepackage{lmodern}
\usepackage{microtype}
\usepackage{amsmath,amssymb,mathtools}
\usepackage{amsthm}
\usepackage{booktabs}
\usepackage{array}
\usepackage{multirow}
\usepackage{xcolor}
\usepackage{enumitem}
\usepackage{needspace}
\usepackage{tikz}
\usetikzlibrary{arrows.meta,positioning,fit}
\usepackage[hidelinks]{hyperref}
\usepackage[nameinlink,noabbrev]{cleveref}
\usepackage{listings}
\usepackage{url}
\usepackage[backend=biber,style=numeric,sorting=nyt,maxbibnames=99]{biblatex}
\addbibresource{references.bib}

\newtheorem{theorem}{Theorem}
\newtheorem{lemma}[theorem]{Lemma}
\newtheorem{proposition}[theorem]{Proposition}
\newtheorem{corollary}[theorem]{Corollary}

\definecolor{deepblue}{RGB}{23,69,124}
\definecolor{darkred}{RGB}{145,32,32}
\definecolor{lightgray}{RGB}{245,246,248}
\ifdefined\ANONYMOUS
  \newcommand{\pdfpaperauthor}{Anonymous}
\else
  \newcommand{\pdfpaperauthor}{Justin Thaler, Kai Zhang, Scott Kominers, Quang Dao}
\fi
\hypersetup{
  colorlinks=true,
  linkcolor=deepblue,
  citecolor=deepblue,
  urlcolor=deepblue,
  pdftitle={Symmetry-Quotient Forgery Attacks on Odd-Characteristic SNOVA},
  pdfauthor={\pdfpaperauthor},
  pdfsubject={Cryptanalysis of the public odd-characteristic SNOVA drafts and previews},
  pdfkeywords={SNOVA, multivariate signatures, forgery attack, NIST PQC}
}
\setlist{nosep,leftmargin=*}
\lstset{basicstyle=\ttfamily\small,backgroundcolor=\color{lightgray},frame=single,breaklines=true}

\newcommand{\F}{\mathbb F}
\newcommand{\Mat}{\operatorname{Mat}}
\newcommand{\Sym}{\operatorname{Sym}}
\newcommand{\rank}{\operatorname{rank}}
\newcommand{\Span}{\operatorname{span}}
\newcommand{\MQ}{\mathsf{MQ}}
\newcommand{\Hashimoto}{\mathsf{Hashimoto}}
\newcommand{\transpose}{\mathsf T}
\newcommand{\SNOVA}{\textsf{SNOVA}}
\newcommand{\QRUOV}{\textsf{QR-UOV}}
\newcommand{\MAYO}{\textsf{MAYO}}
\newcommand{\bits}{\text{ bits}}
\newcommand{\attacktitle}{Symmetry-Quotient Forgery Attacks on Odd-Characteristic \texorpdfstring{\SNOVA}{SNOVA}}

\title{\attacktitle}

\ifdefined\ANONYMOUS
  \author{Anonymous submission}
\else
  \author{%
    Justin Thaler\thanks{a16z crypto research and Georgetown University} \and
    Kai Zhang\thanks{MIT} \and
    Scott Kominers\thanks{Harvard University and a16z crypto research} \and
    Quang Dao\thanks{Carnegie Mellon University and LayerZero Labs}%
  }
\fi
\date{Draft of 24 July 2026}

\begin{document}
\setlength{\emergencystretch}{2em}
\raggedbottom
\maketitle

\ifdefined\DISCLOSURE
\begin{center}
\fcolorbox{darkred}{white}{\parbox{0.91\textwidth}{\centering\small
\textbf{Confidential coordinated disclosure.} This draft is intended for the \SNOVA{} submission team and NIST. Please do not redistribute before the coordinated public-disclosure date.}}
\end{center}
\fi

\begin{abstract}
We give a forgery attack on the odd-characteristic form of \SNOVA{}, a third-round candidate in NIST's Additional Digital Signatures process.  \SNOVA{}'s original characteristic-two parameters have already been broken, and the team's public repair moves to odd characteristic, where the public-key matrices become symmetric; we show that this repair also falls below its targeted security.  In \SNOVA{}'s own cost model, our estimates fall to between $2^{130.3}$ and $2^{138.9}$ Boolean gates at Level~I, about $2^{184.3}$ at Level~III, and between $2^{228.0}$ and $2^{238.8}$ at Level~V, against targets $2^{143}$, $2^{207}$, and $2^{272}$.  All nine Version~2.3 draft parameter sets fall short, by 4.06 to 44.05 bits, as do the six Version~2.4 analysis-preview shapes, by as little as 1.43 bits if they keep the symmetric construction.

The attack builds on the existing Beullens forgery framework (Eurocrypt 2025), which recasts forgery as solving a system of quadratic equations whose hardness is governed by the rank of an associated linear map.  We show that the new symmetry forces that map to depend only on unordered power pairs.  Its universal rank ceiling is therefore $m_1\binom{\ell+1}{2}$ rather than $m_1\ell^2$ (here $\ell$ is the block size and $m_1$ a scheme parameter).  Across the nine Version~2.3 sets, the resulting residual systems range from 48 equations in 112 variables to 96 equations in 224 variables.

The symmetric-square factorization and rank ceiling are exact for every key with symmetric public matrices.  The residual-system dimensions we tabulate additionally assume full quotient rank $\rho=K$ and full affine-constraint rank $\lambda_R=M-K$; we verify these end-to-end on the audited $(28,5,4,4)$ Level-I shape and state them as explicit, checkable genericity assumptions for the other shapes.  We also reconstruct the audited Level-I reduced system from the official reference implementation's own key material.

All bit-security figures inherit the semi-regular and sparse-linear-algebra assumptions used in the \SNOVA{} team's analysis; we do not claim a full-size forgery.  This exact quotient does not apply to the unsymmetrized $q=16$ public matrices.
\end{abstract}

\section{Introduction}\label{sec:intro}

\SNOVA{} is a structured Oil-and-Vinegar signature scheme designed to combine very small public keys with short signatures and fast verification.  NIST advanced \SNOVA{} to the third round of the Additional Digital Signatures process in May~2026, while allowing the surviving candidates to submit further specifications and implementations~\cite{NISTIR8610,NISTRound3}.  The scheme has not, however, had a single stable public parameter generation, and no definitive Round~3 parameter package was public as of 24~July~2026.

The NIST-posted Round~2 package uses $\F_{16}$ and nonsymmetrized public base matrices.  That line has already been heavily damaged: Beullens reduced the security margins by exploiting the structured whipping layer~\cite{Beullens2024SNOVA}, and the subsequent structure-aware wedge attack of Bros et al.\ reports that all Round~2 parameters are broken, while noting a remaining rank conjecture in its complexity analysis~\cite{BrosEtAl2026SNOVAWedge}.  The \SNOVA{} team's later public drafts pursue odd-characteristic variants of the scheme intended to evade these attacks.  Version~2.1 added prime-field variants with symmetric public matrices; Version~2.3 supplies $q=19$ implementations and KATs; the public analysis repository later posted a six-set ``Version~2.4'' preview; and the authors' reformulation paper presents odd characteristic as a principal route for redesigning the scheme~\cite{SNOVA23,SNOVA24Preview,DingEtAl2026Reformulating}.  An odd-characteristic attack therefore targets the public repair strategy rather than the already-vulnerable submitted parameters.

The Version~2.3 attack chapter retains a forgery analysis in the style of Beullens~\cite{Beullens2024SNOVA}.  There, forgery is recast as solving a system of quadratic equations, and the size of that system is taken to be $m_1\ell^2$: the analysis lets the powers $S^0,\dots,S^{\ell-1}$ of the public matrix $S$ contribute all $\ell^2$ ordered pairs $(a,b)$, and it models the resulting equations as essentially random of that size.  The draft itself states that this chapter is inherited from Version~2.0 and still needs updating~\cite{SNOVA23}.

This paper shows that the ordered model is incompatible with the new public symmetry.  Put $A=\F_q[S]$ and $d=\dim_{\F_q}A$; the current irreducible degree-$\ell$ construction has $d=\ell$ and $A\cong\F_{q^\ell}$.  Since $S=S^\transpose$, every element of $A$ is symmetric.  For each public base form $P_i$, exchanging the two labels in
\[
 (\xi,\eta)\longmapsto
 u^\transpose(I_n\otimes\xi)P_i(I_n\otimes\eta)u
\]
does not change the resulting quadratic polynomial in $u$.  Equivalently, the ordered label map $A\otimes_{\F_q}A$ annihilates every tensor $\xi\otimes\eta-\eta\otimes\xi$ and factors through the quotient $\Sym^2_{\F_q}(A)$.  In odd characteristic the alternating tensors form a canonical $\binom d2$-dimensional subspace of this kernel, giving
\[
 A\otimes_{\F_q}A
 \cong \Sym^2_{\F_q}(A)\oplus\bigwedge\nolimits^2_{\F_q}(A).
\]
The alternating subspace is a universal kernel contribution; additional key- or relation-dependent kernel may exist.  Thus the universal homogeneous-rank ceiling is
\[
 K_d=m_1\binom{d+1}{2},
\]
and for the current parameters this is $K=m_1\binom{\ell+1}{2}$ rather than $m_1\ell^2$.  The cap holds for every key generated with symmetric public matrices; it is not a rare MinRank event and cannot be removed by adding more $A,B,Q$ terms (the ``emulsifier'' matrices of the whipping layer, which randomize the public map).

\subsection{Results}

Recall the parameters: $v$ vinegar and $o$ oil matrix variables (so $n=v+o$), block size $\ell$, and $r$ columns per signature block, with $m_1=\lceil or/\ell\rceil$ scalar-expanded public base matrices and $M=or\ell$ scalar verification equations.  Our main algebraic result is the following.

\Needspace{19\baselineskip}
\begin{theorem}[Universal symmetry quotient, informal]\label{thm:intro}
Consider a \SNOVA{} public key for which the power matrix $S$ and the scalar-expanded public base matrices $P_1,\ldots,P_{m_1}$ are symmetric, and write $d=\dim_{\F_q}\F_q[S]$.  After any affine-column substitution in the Beullens framework, the homogeneous quadratic part factors through
\[
 \bigoplus_{i=1}^{m_1}\Sym^2_{\F_q}(\F_q[S])
\]
and has rank at most $K_d=m_1\binom{d+1}{2}$.  For the current irreducible degree-$\ell$ construction, $d=\ell$ and $K_d=K=m_1\binom{\ell+1}{2}$.

For a fixed relation and offsets, let $\rho\le K$ be the actual quotient rank and let $\lambda_R$ be the rank of the induced affine constraints.  Whenever those constraints are consistent with the target, forging is exactly reducible to $\rho$ equations of degree at most two in $\ell(v+o)-\lambda_R$ variables.  In the full-rank case $\rho=K$ and $\lambda_R=M-K$, every target is consistent and the residual dimensions are
\[
 K\text{ equations in }N_{\rm res}=\ell(v+o)-\left(or\ell-K\right)\text{ variables}.
\]
\end{theorem}

The full theorem, including nongeneric ranks, appears as \Cref{thm:rankbound,thm:residual}.  The rank cap comes from the domain of the quadratic feature map.  The $A,B,Q$ emulsifier matrices may randomize the induced map on the symmetric quotient, but they cannot recover the alternating coordinates discarded before emulsification.  For a simple relation, we certify full quotient rank for all nine Version~2.3 shapes.  For the rejection-bypassing relation used in the end-to-end attack, we certify both quotient and affine ranks on the audited $(28,5,4,4)$ Level-I KAT-derived key and two further generated keys of the same shape.  For every other shape, the simultaneous rank conditions for a rejection-bypassing relation are stated explicitly as assumptions in the cost tables (\Cref{sec:validation}).

\paragraph{Scope of exactness.}
The result has three layers with different status.  The symmetric-square factorization and rank ceiling are universal algebraic statements.  The displayed selected-target $K$-in-$N_{\rm res}$ systems require the instance-checkable rank equalities $\rho=K$ and $\lambda_R=M-K$.  Converting those dimensions into gate exponents additionally uses the submission's heuristic semi-regular solver model.  The tables below state the latter two conditions rather than folding them into the universal theorem.

Applying the submission's own MQ methodology gives the Version~2.3 estimates in \Cref{tab:main}.  The submission's estimator admits two conventions for the semi-regular series, differing by whether one homogenizing variable is inserted.  The audited analysis snapshot and the values printed in Version~2.3 Table~11 use the convention without that variable ($p_1=0$); inserting the extra variable ($p_1=1$) yields a stricter, higher estimate.  We adopt the stricter $p_1=1$ convention for the headline figures, so our reported costs are conservative relative to the submission's own numbers; even so, every $q=19$ set falls below target.  Applying the identical quotient to the Version~2.4 analysis preview gives the additional estimates in \Cref{tab:snova24}.

\begin{table}[t]
\centering
\caption{Symmetry-reduced residual systems and estimated attack costs in the full-rank case $\rho=K$, $\lambda_R=M-K$.  A simple relation has certified full quotient rank on all nine shapes.  Simultaneous quotient and affine rank for the rejection-bypassing attack relation is certified end-to-end only on the audited $(28,5,4,4)$ Level-I shape and is an explicit genericity assumption for the other rows.  ``Old'' is the Beullens-forgery estimate in \SNOVA{} Version~2.3 Table~11.  ``New'' uses the stricter $p_1=1$ convention.  Negative margins mean that the estimate is below the category target.}\label{tab:main}
\small
\setlength{\tabcolsep}{4.2pt}
\begin{tabular}{@{}c l c c c r r@{}}
\toprule
Level & Parameters $(v,o,\ell,r)$ & Residual MQ & Old & New & Target & Margin\\
\midrule
I   & $(28,5,4,4)$   & $50$ in $102$ & $194$ & $138.94$ & $143$ & $-4.06$\\
I   & $(48,16,2,2)$  & $48$ in $112$ & $164$ & $130.30$ & $143$ & $-12.70$\\
I   & $(28,4,4,5)$   & $50$ in $98$  & $194$ & $138.94$ & $143$ & $-4.06$\\
\midrule
III & $(40,7,4,4)$   & $70$ in $146$ & $264$ & $184.29$ & $207$ & $-22.71$\\
III & $(72,24,2,2)$  & $72$ in $168$ & $234$ & $184.29$ & $207$ & $-22.71$\\
III & $(38,5,4,5)$   & $70$ in $142$ & $239$ & $184.29$ & $207$ & $-22.71$\\
\midrule
V   & $(50,9,4,4)$   & $90$ in $182$ & $333$ & $227.95$ & $272$ & $-44.05$\\
V   & $(96,32,2,2)$  & $96$ in $224$ & $303$ & $238.77$ & $272$ & $-33.23$\\
V   & $(52,6,4,6)$   & $90$ in $178$ & $333$ & $227.95$ & $272$ & $-44.05$\\
\bottomrule
\end{tabular}
\end{table}

\SNOVA{}'s verifier rejects any signature whose blocks are symmetric.  Because the attack is built from symmetric column relations, one might hope this rule blocks it.  We show it does not: in the square variants the attacker adds fixed offsets that force a nonzero antisymmetric entry into every block, so the forged signature passes the check while the system to be solved is unchanged, and in the rectangular variants the blocks are not symmetric objects in the first place.

We also report an implementation audit against the reference code.  Reconstructing the official Level-I implementation reproduces the reduced system exactly from its shipped key material.  Independently, after fixing one signature column whose cross-term coefficient map is surjective and choosing a target-matching coset of its kernel, the verifier restricts to a sound system of 50 equations of degree at most two in 52 variables (\Cref{sec:crossnormal,sec:validation}).  In the reported structural tests, no shared solution locus or exploitable low-rank shortcut was observed beyond this reduction; this is negative evidence rather than a completeness result.

\subsection{Organization}

\Cref{sec:related} surveys prior work.  \Cref{sec:prelim} fixes notation and recalls the affine-column framework.  \Cref{sec:symmetry} proves the quotient and its multibase Gram extension.  \Cref{sec:attack} derives the smaller system of equations a forger must solve, and \Cref{sec:rejection} shows that the verifier's rule against symmetric signatures does not block the attack.  \Cref{sec:crossnormal} gives the cross-column 50-in-52 normal form.  \Cref{sec:complexity} evaluates our attack's runtime against both the Version~2.3 parameter sets and the Version~2.4 preview.  \Cref{sec:validation} reports the experiments confirming the reduction against \SNOVA{}'s reference implementation and testing its solving-degree model.  \Cref{sec:countermeasures} gives countermeasures.  \Cref{app:noAshortcut} rules out an obvious alternative attack---a MinRank shortcut over the extension field---and \Cref{app:multibase} develops the multibase solver extension.

\section{Related Work}\label{sec:related}

\paragraph{Beullens's forgery attack.}
Beullens identified \SNOVA{}'s structured whipping layer, expressed the public map through ordered bilinear power-pair coordinates, and showed that low-rank combinations of the corresponding emulsifier matrices lead to forgery attacks~\cite{Beullens2024SNOVA}.  The Version~2.3 countermeasure increases the number and diversity of $A,B,Q$ terms to make the induced ordered-coordinate matrix resemble a random matrix.

Our attack uses the same affine-column architecture but a different source of rank loss.  We do not find a low-rank element in a family of otherwise full-rank ordered-coordinate matrices.  We show that the odd-characteristic homogeneous feature vector itself lies in a $K$-dimensional symmetric subspace.  Thus the quotient applies to every key and every relation, even when the induced map is maximally random and full rank on that quotient.

\paragraph{Stable-ideal attacks.}
Cabarcas et al. exploit the action of the commutative group generated by $S$ to design a polynomial-system solver specialized to \SNOVA{}~\cite{Cabarcas2025SNOVA}.  Their work can be viewed as improving the residual-system solving stage.  The present result instead reduces the number of residual equations before solving begins.  The approaches are compatible, but we do not combine their cost benefits here.

\paragraph{Characteristic-two wedge attacks and the odd-characteristic redesign.}
The structure-aware wedge attack exploits \SNOVA{}'s block-ring structure and reports that the NIST Round~2 $q=16$ parameters are broken, while explicitly retaining a conjectural rank component in the analysis~\cite{BrosEtAl2026SNOVAWedge}.  The \SNOVA{} authors subsequently investigated the wedge map and proposed a reformulation emphasizing odd-characteristic and symmetric-form variants~\cite{TsengEtAl2026Wedge,DingEtAl2026Reformulating}.  Their earlier Round~2 security paper explicitly noted that symmetric public matrices over odd-prime fields reduce the number of alternating forms by roughly a factor of two, cautioned that symmetrization may create new attacks, and proposed excluding symmetric signature matrices~\cite{WangEtAl2025Round2Security}.  That observation was qualitative and made in the wedge-attack setting.  Our attack is complementary: it does not apply to the unsymmetrized $q=16$ public matrices, but it attacks the public odd-characteristic repair direction through an exact affine-column quotient that is independent of the wedge-rank conjecture.

\paragraph{Key-recovery attacks.}
Several works interpret the scalar expansion of \SNOVA{} as UOV, or lift it to smaller extension-field UOV instances, and transfer key-recovery attacks~\cite{IkematsuAkiyama2024,LiDing2024,NakamuraEtAl2025}.  Our result is a universal forgery reduction and does not recover the secret oil space.

\paragraph{Relation to the symmetric-algebra key-recovery attack.}
The decomposition
$V\otimes V=\Sym^2(V)\oplus\bigwedge^2(V)$ and the vanishing of the alternating
part under quadratic evaluation are classical; we claim no novelty for that
algebra~\cite{Lam2005}.  The closest cryptanalytic precedent is the
$p$-reduced symmetric-algebra framework of Jin, Pan, He, Gong, and
Ding~\cite{JinEtAl2026}, which generalizes Ran's characteristic-two wedge
attack~\cite{Ran2026} to odd characteristic.  Three differences are worth
stating.  First, theirs is a \emph{key-recovery} XL attack that reconstructs the
oil space, whereas ours is a universal forgery that recovers no secret.  Second,
their algebra is the quotient $\F_q[x_1,\dots,x_n]/(x_1^p,\dots,x_n^p)$ on all
$n=v+o$ variables, in which the polar form is a degree-two element and the
Macaulay computation runs to degree up to $o(p-1)$; our symmetric square is the
degree-two part of the $\ell$-dimensional power algebra $\F_q[S]$, and collapses
only the ordered power pairs of the Beullens feature map.  Third, their cost is
strongly characteristic-sensitive---it grows with $p$ and is ineffective at the
large characteristics of \QRUOV{}---whereas our quotient dimension
$m_1\binom{\ell+1}{2}$ is independent of $p$ for odd $p$.  The two attacks are
therefore complementary, and ours applies directly in the moderate-characteristic
regime $q=19$ that the odd-characteristic \SNOVA{} drafts adopt.  Their own
concrete recommended-parameter evaluation targets \QRUOV{} and reports no
improvement over the existing reconciliation attack at its large characteristic;
they provide no \SNOVA{} cost tables, so no numerical head-to-head on the $q=19$
Version~2.3 sets is available.

\paragraph{Our contribution.}
The contribution is to turn the prior qualitative symmetry warning into the exact
$\Sym^2(A)$ factorization of the Beullens affine-column feature map; derive the
rank and selected-target residual-system consequences for square and rectangular
variants; extend the factorization to the multibase Gram setting; prove
constructively that the proposed signature rejection does not block the attack;
and validate the reduction against the audited Version~2.3 implementation
snapshot and cost methodology.  We found no exact affine-column quotient,
selected-target reduction, or revised cost table in the Version~2.3 attack
chapter or the Version~2.4 analysis preview.  The former continues to use
ordered $\ell^2$ coordinates and explicitly states that it still requires updating.

\section{Preliminaries and the Existing Forgery Framework}\label{sec:prelim}

\subsection{Parameters and scalar expansion}

Let $q$ be a prime power, let $\ell$ be the size of the square matrix blocks, and let $r$ be the number of columns of each signature block.  The public map has $n=v+o$ matrix variables.  Version~2.3 defines
\[
 m_1=\left\lceil\frac{or}{\ell}\right\rceil
 \quad\text{and}\quad
 M=or\ell,
\]
where $m_1$ is the number of scalar-expanded public base matrices and $M$ is the number of scalar output equations.  A signature is an $n$-tuple of blocks in $\Mat_{\ell\times r}(\F_q)$.  Equivalently, writing its columns as
\[
 U=[u_0\mid u_1\mid\cdots\mid u_{r-1}],
 \qquad u_j\in\F_q^{n\ell},
\]
the attack starts with $n\ell$ variables after expressing the other columns affinely in terms of $u_0$.

Let $S\in\Mat_{\ell}(\F_q)$ be the public symmetric matrix with irreducible characteristic polynomial, and put
\[
 \Sigma_a=I_n\otimes S^a\in\Mat_{n\ell}(\F_q),
 \qquad 0\le a<\ell.
\]
For each base matrix $P_i\in\Mat_{n\ell}(\F_q)$ define the bilinear power-pair forms
\begin{equation}\label{eq:B}
 B_i^{(a,b)}(x,y)=x^\transpose\Sigma_aP_i\Sigma_b y.
\end{equation}
The public odd-characteristic algorithms generate each scalar-expanded public matrix symmetrically: diagonal blocks are symmetric and opposite off-diagonal blocks are transposes.  Thus $P_i^\transpose=P_i$ for the public $q=19$ drafts, which we confirm directly against the official Level-I KAT in \Cref{sec:validation}.

\subsection{The Beullens affine-column substitution}

Beullens observed that the \SNOVA{} public map can be written as a linear transformation of the values $B_i^{(a,b)}(u_j,u_k)$~\cite{Beullens2024SNOVA}.  The attacker selects matrices $R_j\in\F_q[S]$ and offsets $v_j\in\F_q^{n\ell}$, with $R_0=I$, and substitutes
\begin{equation}\label{eq:relation}
 u_j=(I_n\otimes R_j)u+v_j,
 \qquad 0\le j<r.
\end{equation}

after identifying $u=u_0$ and taking $v_0=0$ if desired.

For a fixed relation $R=(R_0,\ldots,R_{r-1})$, the substituted verification map for a target $t\in\F_q^M$ can be written
\begin{equation}\label{eq:affine-system}
 E_R z_{\rm ord}(u)+L_Ru+c_R=t.
\end{equation}
Here
\[
 z_{\rm ord}(u)
 =\left(B_i^{(a,b)}(u,u)\right)_{1\le i\le m_1,\ 0\le a,b<\ell}
 \in\F_q^{m_1\ell^2},
\]
$E_R\in\F_q^{M\times m_1\ell^2}$ depends only on public constants and the chosen column relation, while $L_Ru+c_R$ collects all cross and constant terms involving the offsets.

In the ordered-coordinate analysis, if $\rank(E_R)=\rho$, Gaussian elimination supplies $M-\rho$ affine-linear equations and leaves at most $\rho$ independent quadratic equations.  The Version~2.3 attack chapter models $E_R$ as essentially random and bases its concrete estimate on $\rho$ near the ordered-coordinate maximum.  Our observation is that, whenever the public base matrices are symmetric, $E_R$ is applied to a feature vector with deterministic coordinate equalities.

\section{The Symmetry Quotient}\label{sec:symmetry}

\begin{lemma}[Power-pair symmetry]\label{lem:powerpair}
Suppose $S^\transpose=S$ and $P_i^\transpose=P_i$.  Then for every $u\in\F_q^{n\ell}$ and every $0\le a,b<\ell$,
\[
 B_i^{(a,b)}(u,u)=B_i^{(b,a)}(u,u).
\]
\end{lemma}

\begin{proof}
The quantity is a scalar, so it equals its transpose:
\[
 \begin{aligned}
 u^\transpose\Sigma_aP_i\Sigma_bu
 &=\left(u^\transpose\Sigma_aP_i\Sigma_bu\right)^\transpose\\
 &=u^\transpose\Sigma_b^\transpose P_i^\transpose\Sigma_a^\transpose u
 =u^\transpose\Sigma_bP_i\Sigma_au.
 \end{aligned}
\]
We used that every power of the symmetric matrix $S$ is symmetric.
\end{proof}

\begin{theorem}[Structural symmetric-square factorization]\label{thm:structural-sym2}
Let $A=\F_q[S]\subseteq\Mat_\ell(\F_q)$, put $d=\dim_{\F_q}A$, and assume $S^\transpose=S$ and $P_i^\transpose=P_i$.  Let $V=\F_q^{n\ell}$ and let $\mathcal Q_2(V)$ denote the space of formal homogeneous quadratic polynomials on $V$.  Define the linear label map
\[
 \Psi_i:A\otimes_{\F_q}A\longrightarrow\mathcal Q_2(V),
 \qquad
 \Psi_i(\xi\otimes\eta)(u)
 =u^\transpose(I_n\otimes\xi)P_i(I_n\otimes\eta)u.
\]
Then
\[
 \Psi_i(\xi\otimes\eta)=\Psi_i(\eta\otimes\xi)
 \qquad(\xi,\eta\in A).
\]
Consequently $\Psi_i$ annihilates every alternating relation
$\xi\otimes\eta-\eta\otimes\xi$ and factors uniquely through the quotient
$\Sym^2_{\F_q}(A)$.  The combined label map for
$P_1,\ldots,P_{m_1}$ therefore factors through
\[
 \bigoplus_{i=1}^{m_1}\Sym^2_{\F_q}(A),
\]
a space of dimension $m_1\binom{d+1}{2}$.

If the characteristic is odd, antisymmetrization canonically embeds
$\bigwedge^2_{\F_q}(A)$ into $A\otimes_{\F_q}A$, and
\[
 \bigoplus_{i=1}^{m_1}\bigwedge\nolimits^2_{\F_q}(A)
 \subseteq \ker\!\left(\bigoplus_i\Psi_i\right).
\]
This is a universal kernel contribution of dimension
$m_1\binom d2$; equality with the full kernel is not asserted.  If the
minimal polynomial of $S$ has degree $\ell$---in particular, for the
irreducible degree-$\ell$ \SNOVA{} construction---then $d=\ell$.
\end{theorem}

\begin{proof}
Every element of $A$ is a polynomial in the symmetric matrix $S$ and is
therefore symmetric; irreducibility is not needed for this step.  The displayed
map is linear on $A\otimes A$.  For $\xi,\eta\in A$ and $u\in V$, its value is
a scalar and hence equals its transpose:
\[
\begin{aligned}
 \Psi_i(\eta\otimes\xi)(u)
 &=u^\transpose(I_n\otimes\eta)P_i(I_n\otimes\xi)u\\
 &=\left(u^\transpose(I_n\otimes\eta)P_i(I_n\otimes\xi)u\right)^\transpose\\
 &=u^\transpose(I_n\otimes\xi)P_i(I_n\otimes\eta)u
 =\Psi_i(\xi\otimes\eta)(u).
\end{aligned}
\]
Thus the alternating relations lie in the kernel, and the universal property
of the symmetric-square quotient gives the factorization.  Its dimension is
$\binom{d+1}{2}$ per base form.  When $2$ is invertible, the symmetrizer and
antisymmetrizer give the canonical direct decomposition
$A\otimes A\cong\Sym^2(A)\oplus\bigwedge^2(A)$, proving the kernel inclusion.
Finally, $d$ is the degree of the minimal polynomial of $S$; an irreducible
characteristic polynomial of degree $\ell$ therefore gives $d=\ell$.
\end{proof}

Concretely, for $\ell=4$ and $d=\ell$, the ordered label space has
dimension 16, while its symmetric-square quotient has dimension 10.  The six
alternating directions are always killed, although additional kernel may occur
after the public output map is applied.  With $m_1=5$, the universal quotient
removes 30 ordered coordinates; in the audited Level-I relation the induced
quotient map has full rank 50, giving the observed 30-dimensional left kernel.

Define the unordered feature vector
\[
 z_{\rm sym}(u)
 =\left(B_i^{(a,b)}(u,u)\right)_{1\le i\le m_1,\ 0\le a\le b<\ell}
 \in\F_q^K,
 \qquad K=m_1\binom{\ell+1}{2}.
\]
Let
\[
 D:\F_q^K\longrightarrow\F_q^{m_1\ell^2}
\]
be the duplication map that sends an unordered coordinate $(i,\{a,b\})$ to both ordered positions $(i,a,b)$ and $(i,b,a)$ when $a\ne b$, and to the single diagonal position when $a=b$.  Then
\begin{equation}\label{eq:dup}
 z_{\rm ord}(u)=Dz_{\rm sym}(u).
\end{equation}

\begin{figure}[t]
\centering
\resizebox{\textwidth}{!}{%
\begin{tikzpicture}[node distance=11mm,>=Latex,
 box/.style={draw,rounded corners,align=center,minimum height=10mm,minimum width=29mm,fill=lightgray}]
\node[box] (u) {$u\in\F_q^{n\ell}$};
\node[box,right=of u] (sym) {$z_{\rm sym}(u)$\\$K=m_1\binom{\ell+1}{2}$};
\node[box,right=of sym] (ord) {$z_{\rm ord}(u)$\\$m_1\ell^2$ positions};
\node[box,right=of ord] (out) {quadratic output\\in $\F_q^M$};
\draw[->] (u)--node[above]{quadratic}(sym);
\draw[->] (sym)--node[above]{$D$}(ord);
\draw[->] (ord)--node[above]{$E_R$}(out);
\draw[->,deepblue,thick,bend right=23] (sym) to node[below]{$\overline E_R=E_RD$} (out);
\end{tikzpicture}%
}
\caption{The ordered feature vector lies in the image of a $K$-dimensional duplication map.  Randomizing $E_R$ cannot increase the rank of the composite beyond $K$.}\label{fig:factor}
\end{figure}

\begin{theorem}[Symmetry rank bound]\label{thm:rankbound}
For every public key whose power matrix and scalar-expanded public base matrices are symmetric, and for every affine-column relation $R$, the following holds.  The homogeneous quadratic part in \Cref{eq:affine-system} factors as
\[
 \overline E_R z_{\rm sym}(u),
 \qquad \overline E_R=E_RD\in\F_q^{M\times K}.
\]
In particular,
\[
 \rank(\overline E_R)\le K=m_1\binom{\ell+1}{2}.
\]
\end{theorem}

\begin{proof}
The structural factorization is \Cref{thm:structural-sym2}.  In the power
basis $1,S,\ldots,S^{\ell-1}$, it becomes the coordinate identity
\Cref{eq:dup}.  Substituting that identity into \Cref{eq:affine-system}
gives the displayed factorization, and the rank bound follows from the number
of columns of $\overline E_R$.
\end{proof}

\paragraph{Why the emulsifier does not remove the relation.}
The Version~2.3 changes use many independently generated $A_{i,\alpha},B_{i,\alpha},Q_{i,\alpha,1},Q_{i,\alpha,2}$ matrices so that the ordered-coordinate $E_R$ resembles a general matrix rather than the repeated block-diagonal matrix attacked in the first-round construction.  This may make $\overline E_R$ a generic map on its $K$-dimensional domain.  It cannot make that domain larger.  The lost antisymmetric coordinates vanish before the emulsifier is applied.

\paragraph{The reduction is universal, not a weak-key event.}
Prior analyses search over relations $R$ to find an unusually low rank.  Here the cap $K$ holds for every relation and every public key satisfying the public-matrix symmetry rule.  A further quotient-rank drop below $K$, when accompanied by full affine rank, gives no larger residual size parameters and lowers the candidate estimator used below.  Because quotient and affine ranks can interact, both should nevertheless be audited on each concrete relation.

\subsection{The Gram viewpoint and several base vectors}\label{sec:gram}

The same argument gives a clean generalization of the attack.  Choose
$c$ independent base vectors $x_1,\ldots,x_c$ and an $A$-valued linear
column relation
\begin{equation}\label{eq:cbase}
 u_j=\sum_{s=1}^{c}(I_n\otimes R_{j,s})x_s+v_j,
 \qquad R_{j,s}\in A.
\end{equation}
Put $V_c=\F_q^c\otimes_{\F_q}A$, of dimension $cd$, where
$d=\dim_{\F_q}A$.  For each public base matrix $P_i$, the homogeneous
restriction of the verifier is a symmetric Gram pairing on $V_c$.  We therefore
obtain the following exact extension.

\begin{proposition}[$c$-base Gram factorization]\label{prop:cbase}
Under \Cref{eq:cbase}, the homogeneous feature map factors through
\[
 \bigoplus_{i=1}^{m_1}\Sym^2_{\F_q}(V_c),
\]
and hence through at most
\[
 K_{c,d}=m_1\binom{cd+1}{2}
\]
quadratic coordinates.  For the current irreducible degree-$\ell$
construction, $d=\ell$ and this becomes
$K_c=m_1\binom{c\ell+1}{2}$.  If $c=r$ and the $r\times r$ relation matrix over
$A$ is invertible, the relation is merely a public change of basis among the
signature columns and imposes no restriction on the signature space.
\end{proposition}

\begin{proof}
Apply the proof of \Cref{thm:structural-sym2} to the labels
$e_s\otimes\xi\in V_c$.  Symmetry of the public form identifies a pair of
labels with its reversal, so the universal feature space is the symmetric
square of $V_c$.  An invertible $A$-linear column change is bijective.
\end{proof}

The Gram formulation also closes a broader search over arbitrary coefficient
matrices in the current irreducible degree-$\ell$ construction.  For a general
relation with $R_{j,s}\in\Mat_\ell(\F_q)$, define the right power-orbit space
\[
 \mathcal W_s=\Span_{\F_q}\{S^aR_{j,s}:0\le a<\ell,\ 0\le j<r\}.
\]
If base $s$ is active, some $R_{j,s}$ is nonzero.  The map
$A\to\Mat_\ell(\F_q)$, $\xi\mapsto\xi R_{j,s}$, is injective because every
nonzero $\xi\in A$ is invertible.  Therefore
\begin{equation}\label{eq:orbit-lower}
 \dim_{\F_q}\mathcal W_s\ge\ell.
\end{equation}
An $A$-valued relation attains equality for each active base.  Thus allowing
arbitrary block coefficients cannot use fewer than $\ell$ power directions per
active base; it normally makes the orbit spaces much larger.  Here dimensions
across different bases are counted in the labelled direct sum
$\bigoplus_{s=1}^c\mathcal W_s$: as unlabelled subspaces of
$\Mat_\ell(\F_q)$, distinct $A$-valued orbit spaces may coincide.  In 500 trials
per $c\in\{1,2,3,4\}$ on the Level-I shape, the labelled sum
$\sum_s\dim\mathcal W_s$ was always $c\ell$ for $A$-valued relations and
$c\ell^2$ for unrestricted random $\ell\times\ell$ coefficients.  Accidental
rank defects of the public Gram map remain possible, but a broader coefficient
search cannot improve the universal feature count by shrinking an active
power-orbit space below the $A$-valued construction.

For the official Level-I constants, the quotient ranks for $c=1,2,3,4$ are
respectively $50,80,80,80$, equal to the generic maximum
$\min\{K_c,80\}$ in every tested relation.  Thus the tested $c\ge2$ linear
feature-to-output maps are surjective.  This removes only the \emph{linear}
question whether a target lies in the image of that map; it does not show that
the nonlinear Gram-feature map realizes a point in every target fibre.  The
remaining task is a structured nonlinear inversion problem.  We discuss the
resulting candidate multibase architecture and its present solver uncertainty in
\Cref{app:multibase}.

\section{The Symmetry-Assisted Forgery Attack}\label{sec:attack}

Let
\[
 \rho=\rank(\overline E_R)\le K.
\]
Choose a full-rank matrix $W_R\in\F_q^{(M-\rho)\times M}$ whose rows span the left kernel of $\overline E_R$.  Multiplying \Cref{eq:affine-system} by $W_R$ eliminates every quadratic term:
\begin{equation}\label{eq:linear}
 W_RL_Ru=W_R(t-c_R).
\end{equation}
Let $\lambda_R=\rank(W_RL_R)$.  Gaussian elimination expresses $\lambda_R$ coordinates of $u$ as affine functions of the remaining
\[
 N_R=n\ell-\lambda_R
\]
coordinates.  Substitution into a row basis of \Cref{eq:affine-system} leaves at most $\rho$ quadratic equations.

\begin{theorem}[Residual MQ reduction]\label{thm:residual}
Fix a relation and offsets.  Let $\rho=\rank(\overline E_R)$, let $W_R$ have
rows spanning the left kernel of $\overline E_R$, and put
$\lambda_R=\rank(W_RL_R)$.  Choose $V_R\in\F_q^{\rho\times M}$ so that
\[
 T_R=\begin{bmatrix}V_R\\W_R\end{bmatrix}\in\operatorname{GL}_M(\F_q).
\]
If the affine system
\[
 W_RL_Ru=W_R(t-c_R)
\]
is consistent, then one can compute a bijective affine parametrization
$u=u_0+Jx$ of its solution set, with
$x\in\F_q^{n\ell-\lambda_R}$.  Under this parametrization, solutions of the
original $M$ verification equations are in bijection with solutions of the
$\rho$ equations
\[
 V_R\overline E_R z_{\rm sym}(u_0+Jx)
 +V_RL_R(u_0+Jx)=V_R(t-c_R),
\]
each of degree at most two.  Thus forging reduces in polynomial time to
solving at most $\rho$ genuinely quadratic equations in
$n\ell-\lambda_R$ variables, where
\[
 \rho\le m_1\binom{\ell+1}{2},
 \qquad
 \lambda_R\le M-\rho.
\]

In the full-rank case $\rho=K$ and $\lambda_R=M-K$, the matrix $W_RL_R$ has
full row rank and the affine system is consistent for every target.  The
residual dimensions are therefore
\begin{equation}\label{eq:resparams}
 \boxed{
 K=m_1\binom{\ell+1}{2}\text{ equations of degree at most two in }
 N_{\rm res}=n\ell-M+K\text{ variables}.}
\end{equation}
Every residual solution lifts by linear algebra to a signature satisfying all
$M$ original verification equations.
\end{theorem}

\begin{proof}
The rows of $W_R$ span the complete left kernel of $\overline E_R$, so
multiplication by the invertible matrix $T_R$ transforms the original system
into two equivalent blocks:
\[
\begin{aligned}
 V_R\overline E_Rz_{\rm sym}(u)+V_RL_Ru&=V_R(t-c_R),\\
 W_RL_Ru&=W_R(t-c_R).
\end{aligned}
\]
Consistency of the second block gives an affine solution space of dimension
$n\ell-\lambda_R$ and a bijective parametrization $u=u_0+Jx$.  Substitution
into the first block gives exactly $\rho$ equations of degree at most two;
some may lose their quadratic part after restriction, so there are at most
$\rho$ genuinely quadratic equations.  Since $T_R$ is invertible and the
parametrization is bijective onto the affine solution space, this operation
preserves and reflects solutions.  In the full-rank case $W_RL_R$ is a
surjective map onto $\F_q^{M-K}$, which proves target consistency and the
displayed dimensions.
\end{proof}

\begin{proposition}[Effect of a rank defect]
\label{prop:affine-rank}
Two departures from the full-rank case behave differently.
\begin{enumerate}
\item A \emph{quotient-rank} defect $\rho<K$ accompanied by full affine rank
$\lambda_R=M-\rho$ leaves $\rho$ equations of degree at most two in
$n\ell-M+\rho$ variables and imposes no condition on the target.  Both the
equation count and the variable count are no larger than in the full-quotient
case, and the estimator used in this paper assigns no larger cost.  We do not
claim a solver-independent monotonicity theorem for arbitrary structured MQ
families.
\item An \emph{affine-rank} defect $\lambda_R<M-\rho$ imposes
$(M-\rho)-\lambda_R$ independent consistency conditions on the target.  For a
uniform target these hold with probability exactly
$q^{-((M-\rho)-\lambda_R)}$.  For the scheme's hash-to-field encoding this is
the corresponding random-oracle approximation, so salt resampling incurs that
inverse probability in expectation unless the attacker finds a relation and
offsets with larger affine rank.
\end{enumerate}
\end{proposition}

\begin{proof}
For (1), full row rank of $W_RL_R$ makes the affine block surjective and leaves
the dimensions stated by \Cref{thm:residual}.  The final cost comparison is a
statement about the explicit estimator evaluated in this paper.  For (2), the
image of $W_RL_R$ is a $\lambda_R$-dimensional subspace of
$\F_q^{M-\rho}$.  Since $W_R$ has full row rank, a uniform $t$ induces a
uniform right-hand side $W_R(t-c_R)$, which lies in that image with the stated
probability.
\end{proof}

\paragraph{The tabulated systems are the generic full-rank case.}
The residual dimensions in our tables are stated for $\rho=K$ and $\lambda_R=M-K$, where the affine subsystem has full row rank, \Cref{eq:linear} is unconditionally solvable, and the reduction is a genuine selected-target forgery.  A simple public relation attains $\rho=K$ on all nine Version~2.3 shapes and across the larger Level-I relation search.  The relation equipped with the fixed skew offsets of \Cref{lem:skew} is audited jointly: both equalities hold end-to-end on the KAT-derived $(28,5,4,4)$ Level-I key and two further generated keys of that same shape (\Cref{sec:validation}).  The first certificate does not imply the second, because quotient and affine ranks depend on the chosen relation and offsets.  Consequently the simultaneous full-rank conditions are directly evidenced only for the audited shape and remain explicit genericity assumptions for every other tabulated shape; a defect would alter the residual dimensions or add the target-consistency cost of \Cref{prop:affine-rank}.

\paragraph{Attack algorithm.}
Full affine rank removes the linear target-membership obstruction, but it does not by itself prove that every restricted nonlinear fiber contains a root.  Under the same residual-solvability, guessing, and specialization assumptions used by the cost model in \Cref{sec:complexity}, the reduction gives a selected-message forgery attack in the key-only setting: the attacker targets a message of its choice and makes no signing queries.  If the affine block is inconsistent, or if a sampled residual fiber has no root, the procedure resamples salts or relations.  For an arbitrary message $\mathcal M$:
\begin{enumerate}
\item sample a salt and compute the target $t=H(\mathcal M\|\textsf{salt})$;
\item choose a public column relation $R$ and offsets $v_j$;
\item construct $\overline E_R,L_R,c_R$ and row-reduce as above;
\item solve the residual underdetermined MQ system;
\item reconstruct $u$ and then the complete signature matrix $U$ using \Cref{eq:relation}.
\end{enumerate}
The first, second, fourth, and fifth steps are the same attack architecture used by Beullens.  The new ingredient is that step~3 must quotient the homogeneous features by symmetry before estimating their rank.

\section{Bypassing the Symmetric-Block Rejection}\label{sec:rejection}

The \SNOVA{} team's 2025 security paper proposed excluding symmetric signature matrices as a mitigation for attacks that odd-characteristic symmetrization might introduce~\cite{WangEtAl2025Round2Security}.  Version~2.3 implements a thresholded version of this rule: for odd $q$, the verifier rejects square signatures with any symmetric block when $\ell>2$, and with more than $\lfloor n/4\rfloor$ symmetric blocks when $\ell=2$~\cite{SNOVA23}.

The affine-column attack has fixed offsets in addition to its free vector $u$.  Those offsets can enforce nonsymmetry identically.

\begin{lemma}[A constant-skew column relation]\label{lem:skew}
Assume $\ell\ge2$ and that the characteristic polynomial of $S$ is
irreducible.  Let $e_0,e_1$ be two distinct standard basis vectors of
$\F_q^\ell$.  There is a unique $R_1\in\F_q[S]$ such that
\[
 e_0^\transpose R_1=e_1^\transpose.
\]
It is computable by one $\ell\times\ell$ linear solve.  With $R_0=I$, choose
offsets for the first two signature columns so that, in each matrix block,
\[
 (v_1)_0-(v_0)_1=\delta
\]
for a fixed $\delta\in\F_q^\times$.  Then every assignment to $u$ satisfies
\[
 U[0,1]-U[1,0]=\delta
\]
in that block.  In particular, the block is never symmetric.
\end{lemma}

\begin{proof}
Consider the linear map
\[
 \theta:\F_q[S]\longrightarrow\F_q^\ell,
 \qquad R\longmapsto e_0^\transpose R.
\]
If $R\ne0$, irreducibility makes $R$ invertible, so
$e_0^\transpose R\ne0$.  Hence $\theta$ is injective; both spaces have
dimension $\ell$, so it is an isomorphism.  This proves existence and
uniqueness.  Constructively, let $G$ be the $\ell\times\ell$ matrix whose
$a$th row is $e_0^\transpose S^a$.  Solve $c^\transpose G=e_1^\transpose$ and
set $R_1=\sum_{a=0}^{\ell-1}c_aS^a$.

Under the relation $u_0=u+v_0$ and
$u_1=(I_n\otimes R_1)u+v_1$, the variable-dependent parts of the entries
$U[0,1]$ and $U[1,0]$ are $e_0^\transpose R_1u$ and
$e_1^\transpose u$, respectively, and therefore coincide.  Their difference
is the chosen offset constant $\delta$.
\end{proof}

The construction can be applied independently to every block.  It does not change the homogeneous quadratic map $\overline E_R$; offsets affect only $L_R$ and $c_R$.  For the square $\ell=2$ sets it makes every block nonsymmetric, which is stronger than the verifier's threshold.  For the square $\ell=4$ sets it satisfies the zero-symmetric-block requirement.  The rectangular variants have blocks in $\Mat_{\ell\times r}$ with $r\ne\ell$, so matrix symmetry is not defined and no such rejection is applied.  This bypass is algebraically universal.  Whether a chosen relation and these offsets simultaneously give full affine rank is a separate, instance-checkable condition; the Level-I end-to-end audit verifies that condition for the relation used there.

\section{A Cross-Column Gram Normal Form}\label{sec:crossnormal}

The multibase Gram viewpoint yields a second exact reduction for the official Level-I square parameters.  Write the four signature columns as $x_0,x_1,x_2,x_3\in\F_{19}^{132}$.  The complete symmetry-reduced feature space consists of four self blocks and six cross blocks.  For a pair $s<t$, let
\[
 G_{st}(x_s,x_t)
 =\bigl(B_i^{(a,b)}(x_s,x_t)\bigr)_{1\le i\le5,\ 0\le a,b<4}
 \in\F_{19}^{80},
\]
and let $E_{st}:\F_{19}^{80}\to\F_{19}^{80}$ be the corresponding public output map.

For the self contribution of column $x_1$, let
\[
 z_1^{\mathrm{self}}(x_1)\in\F_{19}^{50},
 \qquad
 E_{11}^{\mathrm{self}}:\F_{19}^{50}\longrightarrow\F_{19}^{80},
 \qquad
 Q_1(x_1)=E_{11}^{\mathrm{self}}z_1^{\mathrm{self}}(x_1).
\]
Here $z_1^{\mathrm{self}}$ is quadratic, whereas
$E_{11}^{\mathrm{self}}$ is linear.

\begin{proposition}[Official Level-I cross-pair certificate]\label{prop:crosspair}
For the official Version~2.3 Level-I $q=19$,
$(v,o,\ell,r)=(28,5,4,4)$ public constants, all six maps $E_{st}$ have rank
80.  The full $80\times680$ symmetry-reduced feature map, the span of all self
features, and the span of all cross features also each have rank 80.  The
linear self-feature output map $E_{11}^{\mathrm{self}}$ has rank 50.
\end{proposition}

\begin{proof}[Author-reported exact computation]
Each matrix is reconstructed from the official public constants and
row-reduced by exact Gaussian elimination over $\F_{19}$; no floating-point
rank test is used.  The authors' audit scripts report the indicated pivot counts
for all six cross maps and for the combined, self, cross, and
$E_{11}^{\mathrm{self}}$ matrices.
\end{proof}

Fix $x_0,x_2,x_3$.  After absorbing every verifier contribution
independent of $x_1$ into $d\in\F_{19}^{80}$, all remaining cross terms are
linear in $x_1$.  Let
\[
 C_{x_0,x_2,x_3}:\F_{19}^{132}\longrightarrow\F_{19}^{80}
\]
denote their sum.  The verifier equations restricted to these fixed columns
have the exact form
\begin{equation}\label{eq:cross-fiber}
 C_{x_0,x_2,x_3}(x_1)+Q_1(x_1)=d.
\end{equation}
For the clean two-column restriction $x_2=x_3=0$ this specializes to
\[
 C_{x_0,0,0}(x_1)=E_{01}G_{01}(x_0,x_1)=:C_{x_0}(x_1).
\]
The self term must be retained: it is quadratic in $x_1$ and takes values in
the fixed 50-dimensional linear subspace
$\mathcal S_1=\operatorname{im}(E_{11}^{\mathrm{self}})$ of the
output-coordinate space.

\begin{corollary}[A sound 50-in-52 affine slice]\label{cor:cross-slice}
Let $C=C_{x_0,x_2,x_3}$ and assume its publicly computable coefficient matrix
has rank 80.  Choose any $a\in\F_{19}^{132}$ satisfying $C(a)=d$, and let
$J:\F_{19}^{52}\to\ker(C)$ be a linear isomorphism.  Then the solutions of
\Cref{eq:cross-fiber} on the affine slice $a+\ker(C)$ are in bijection with
the solutions of
\begin{equation}\label{eq:cross-50-52}
 Q_1(a+Jw)=0,
 \qquad w\in\F_{19}^{52}.
\end{equation}
After choosing coordinates on $\mathcal S_1$, this is a system of 50 equations
of degree at most two in 52 variables over $\F_{19}$.
\end{corollary}

\begin{proof}
Rank 80 makes $C$ surjective, so $a$ exists and
$\dim\ker(C)=132-80=52$.  For $x_1=a+Jw$, $C(x_1)=d$, and
\Cref{eq:cross-fiber} is therefore equivalent to $Q_1(a+Jw)=0$.  By
\Cref{prop:crosspair}, all values of $Q_1$ lie in the 50-dimensional linear
subspace $\mathcal S_1$, so any linear coordinate isomorphism
$\mathcal S_1\cong\F_{19}^{50}$ yields 50 scalar equations.  The affine
parametrization is bijective on the selected slice, proving soundness and the
claimed correspondence.
\end{proof}

The rank condition on $C$ is not universal.  For the clean two-column slice,
$x_0=0$ gives rank zero.  The rank is publicly testable; on the audited key, all
100 sampled choices of $x_0$ gave $\rank(C_{x_0})=80$.  The attacker can retain
any passing choice and solve the linear system $C(a)=d$.  Fixed columns may also
be chosen to address the verifier's symmetric-block rejection, but then the rank
of the resulting summed cross map
must be tested anew; the clean $C_{x_0}$ certificate alone does not prove this
joint condition.  The resulting slice is sound but not complete: it contains
only the points in one target-matching coset of the kernel.  The 50 equations
in \Cref{eq:cross-50-52} are generally inhomogeneous after the shift by $a$.
Underdetermination by two variables does not prove existence.  Under the stronger independent random-function heuristic, the expected number of roots is $19^2$ and the zero-root probability is $(1-19^{-50})^{19^{52}}\approx e^{-19^2}$.  Random quadratic systems need not satisfy that independence assumption, so this remains only a solvability heuristic for the target-matching slice.

The normal form is useful for both cryptanalysis and auditing because it exposes all remaining nonlinearity in a nearly square system.  It does not by itself lower the conservative headline exponent: the Hashimoto estimator for the original 50-in-102 system already specializes almost all surplus variables before its square solve.  We therefore tested whether the restricted 50-dimensional quadratic space has further public structure.  Across 100 random restricted systems, the common polar radical was always zero, the quadratic span always had dimension 50, and every individual form had rank 51 or 52.  On five independently generated keys, the adjoint/commutant algebra had dimension one.  Random sampling of $10^5$ linear combinations in the analogous Version~2.4 50-in-44 system found minimum rank 42, consistent with a generic high-rank matrix space.  These checks provide negative evidence against an immediate radical, decomposition, or low-rank shortcut, while leaving specialized multigraded solving as an open route.

\section{Concrete Complexity Estimates}\label{sec:complexity}

\subsection{Methodology}

We evaluate the new residual systems with the methodology already used in the \SNOVA{} submission and in Beullens's paper.  For a determined or overdetermined MQ system, the model estimates sparse Wiedemann/FXL from a semi-regular operating degree.  Hashimoto's method reduces an underdetermined system of $m$ equations in $n$ variables to a collection of smaller determined systems, optimizing over its parameters $(a,k)$~\cite{Hashimoto2023,Beullens2024SNOVA}.

The Version~2.3 model evaluates an $R_{\rm old}$-equation residual in
$n\ell-M+R_{\rm old}$ variables, up to its field-to-gate conversion.  Its
ordered-pair analysis takes $R_{\rm old}$ from an $m_1\ell^2$-dimensional
feature space.  In the full-rank symmetry case we instead evaluate
\[
 \Hashimoto(n\ell-M+K,K,q),
\]
with $K=m_1\binom{\ell+1}{2}$ as in \Cref{eq:resparams}.

The estimator admits two conventions for the semi-regular series.  In the
analysis snapshot audited for this manuscript, the values in Version~2.3
Table~11 are numerically reproduced by using no extra homogenizing variable
($p_1=0$).  Inserting one extra variable ($p_1=1$) increases the estimates by
about 4.25 bits in most rows.  We adopt the stricter, author-selected $p_1=1$
sensitivity convention in the headline table---so our figures are conservative
relative to those $p_1=0$ numbers---and give both conventions in
\Cref{tab:both}.  The exact repository commit and estimator file must be fixed
in the public artifact; without that pin, ``current script'' is not a stable
provenance claim.

\subsection{The determined systems actually costed by the estimator}
\label{sec:actual-cores}

The estimator does not run a Gr\"obner-basis algorithm directly on the full
underdetermined residual systems displayed in our parameter tables.  It first
specializes surplus variables and then applies HybridMQ or Hashimoto's
reduction to smaller determined or overdetermined cores.  Under the stricter
$p_1=1$ convention, the load-bearing systems are summarized in
\Cref{tab:actual-cores}.

\begin{table}[t]
\centering
\caption{Dominant determined cores selected by the audited estimator snapshot under the
stricter $p_1=1$ convention.  The degree is the semi-regular operating degree
used in the cost calculation.  Degrees marked $^\dagger$ are Hashimoto blocks
that were not produced as explicit systems; they are the predicted operating
degree only, not F4-audited (\Cref{sec:validation}).}
\label{tab:actual-cores}
\small
\setlength{\tabcolsep}{4pt}
\begin{tabular}{@{}c l c c c r@{}}
\toprule
Level & Family & Public residual & Method & Dominant core & Degree\\
\midrule
I   & $\ell=4$ & $50/102$ or $50/98$ & square specialization & $50/41$ & 14\\
I   & $\ell=2$ & $48/112$ & Hashimoto $(a,k)=(2,11)$ & $46/35$ & $11^\dagger$\\
III & $\ell=4$ & $70/146$ or $70/142$ & square specialization & $70/55$ & 16\\
III & $\ell=2$ & $72/168$ & Hashimoto $(a,k)=(2,15)$ & $70/55$ & $16^\dagger$\\
V   & $\ell=4$ & $90/182$ or $90/178$ & square specialization & $90/73$ & 21\\
V   & $\ell=2$ & $96/224$ & Hashimoto $(a,k)=(2,23)$ & $94/71$ & $18^\dagger$\\
\bottomrule
\end{tabular}
\end{table}

For example, the strict Level-I $\ell=4$ estimate first specializes the $50/102$
system to a square $50/50$ system and then guesses nine more variables, leaving
the $50/41$ overdetermined core whose modeled degree is 14.  A solving-degree
experiment on the untouched $50/102$ presentation would therefore not directly
validate the quoted work factor; the relevant experiment is on these
estimator-selected cores.

\paragraph{Sensitivity to solving-degree error.}
The narrow Level-I $\ell=4$ rows are genuinely sensitive to the semi-regular
prediction.  Raising the dominant degree from 14 to 15 multiplies the squared
Macaulay dimension by approximately $\bigl(\binom{56}{15}/\binom{55}{14}\bigr)^2
=(56/15)^2$, or about $2^{3.8}$.  Thus degree 15 would place the strict estimate
close to, but still approximately below, the $2^{143}$ target; degree 16 would
likely erase the claimed shortfall.  These two rows should therefore be read as
narrow estimated shortfalls.

The larger instances have substantially more tolerance.  One additional degree
in the Level-III and Level-V dominant cores costs only roughly four bits in the
same monomial-count model, compared with margins of 22.71--44.05 bits.  The
Level-I $\ell=2$ row has a 12.70-bit margin.  Accordingly, moderate
anti-regularity could change the two narrow Level-I $\ell=4$ conclusions without
changing the Level-III or Level-V category conclusions.

\subsection{All odd-characteristic parameter sets}

For the nine $q=19$ parameter shapes in Version~2.3, \Cref{eq:resparams} gives the residual systems shown in \Cref{tab:both}.  The category targets are the values used by the submission: 143, 207, and 272 bits for Levels I, III, and V.

\begin{table}[t]
\centering
\caption{Estimates under both semi-regular conventions, conditional on $\rho=K$ and $\lambda_R=M-K$ for the attack relation.  A simple relation has certified full quotient rank on all nine shapes; simultaneous quotient and affine rank for the rejection-bypassing relation is directly certified only on the audited $(28,5,4,4)$ Level-I shape.  The $p_1=1$ column is the stricter convention used in the headline figures.  The $p_1=0$ column numerically matches Version~2.3 Table~11 and the audited analysis snapshot.}\label{tab:both}
\small
\setlength{\tabcolsep}{4pt}
\begin{tabular}{@{}c l c r r r@{}}
\toprule
Level & Parameters & Residual MQ & $p_1=0$ & $p_1=1$ & Target\\
\midrule
I   & $(28,5,4,4)$  & $50/102$ & 134.69 & 138.94 & 143\\
I   & $(48,16,2,2)$ & $48/112$ & 126.05 & 130.30 & 143\\
I   & $(28,4,4,5)$  & $50/98$  & 134.69 & 138.94 & 143\\
III & $(40,7,4,4)$  & $70/146$ & 180.05 & 184.29 & 207\\
III & $(72,24,2,2)$ & $72/168$ & 181.89 & 184.29 & 207\\
III & $(38,5,4,5)$  & $70/142$ & 180.05 & 184.29 & 207\\
V   & $(50,9,4,4)$  & $90/182$ & 223.70 & 227.95 & 272\\
V   & $(96,32,2,2)$ & $96/224$ & 234.52 & 238.77 & 272\\
V   & $(52,6,4,6)$  & $90/178$ & 223.70 & 227.95 & 272\\
\bottomrule
\end{tabular}
\end{table}

The Level-I $\ell=4$ rows have the narrowest margin: 4.06 bits under the stricter convention.  The convention choice moves the estimate by about 4.25 bits, more than this margin, yet both conventions place these rows below target (134.69 and 138.94 bits against 143; \Cref{tab:both}), so the shortfall does not turn on that choice.  We nonetheless describe the Level-I rows as \emph{estimated below target}.  The remaining rows have margins of 12.70 to 44.05 bits under the same convention.

\subsection{Version 2.4 analysis preview}\label{sec:snova24}

The public \SNOVA{} analysis repository was updated with six parameter shapes labeled Version~2.4.  As of 24~July~2026, the official implementation repository's supporting documentation still listed Version~2.3 and no matching Version~2.4 specification or implementation package was public~\cite{SNOVA23,SNOVA24Preview}.  The symmetry quotient is algebraic and does not depend on the numerical parameters.  If the preview retains the symmetric public-matrix construction, the stricter $p_1=1$ convention gives \Cref{tab:snova24}.

\begin{table}[t]
\centering
\caption{Symmetry-reduced estimates for the Version~2.4 analysis preview, conditional on retention of the symmetric construction and on $\rho=K$, $\lambda_R=M-K$.  These are not definitive Round~3 parameters; the table evaluates the public analysis-repository shapes under the same $p_1=1$ convention used for Version~2.3.}\label{tab:snova24}
\small
\setlength{\tabcolsep}{5pt}
\begin{tabular}{@{}c l c r r r@{}}
\toprule
Level & Parameters $(v,o,\ell,r)$ & Residual MQ & New & Target & Margin\\
\midrule
I   & $(26,5,4,4)$   & $50/94$  & 138.94 & 143 & $-4.06$\\
I   & $(43,17,2,2)$  & $51/103$ & 141.57 & 143 & $-1.43$\\
III & $(37,7,4,4)$   & $70/134$ & 184.29 & 207 & $-22.71$\\
III & $(64,25,2,2)$  & $75/153$ & 194.57 & 207 & $-12.43$\\
V   & $(44,6,4,6)$   & $90/146$ & 227.95 & 272 & $-44.05$\\
V   & $(90,33,2,2)$  & $99/213$ & 246.84 & 272 & $-25.16$\\
\bottomrule
\end{tabular}
\end{table}

The first Level-I margin is the same narrow 4.06-bit estimate as in Version~2.3, and the second is only 1.43 bits.  Those two rows should be described as estimated shortfalls.  The Level-III and Level-V margins are substantially larger.  Reparameterizing within the same symmetric construction must therefore use $m_1\binom{\ell+1}{2}$, not $m_1\ell^2$, as the homogeneous rank ceiling.

\subsection{Interpretation of the estimates}

All displayed residual dimensions use the simultaneous full-rank equalities
$\rho=K$ and $\lambda_R=M-K$ for the concrete relation and offsets used by
the attack.  A simple relation attains the first equality on every Version~2.3
shape, but that does not certify the pair for a rejection-bypassing relation.
The pair is certified in the audited $(28,5,4,4)$ Level-I end-to-end instances
and is an explicit genericity assumption elsewhere, including the Version~2.4
preview.  The parameter tables therefore require four attack-specific
conditions: the evaluated public construction satisfies the stated symmetry
hypothesis; the same relation and offsets attain the quoted quotient rank; the
associated affine system has the quoted rank and is target-consistent (automatic
when $\lambda_R=M-\rho$); and the residual fiber is nonempty, with every
guessing or specialization step preserving a solution at the success probability
charged by the estimator.  Conditional on
those points, the gate estimates inherit three standard solver assumptions:
\begin{enumerate}
\item the residual systems behave close enough to semi-regular systems for the operating-degree calculation to be meaningful;
\item sparse linear algebra has the costs assumed by the submission's model;
\item field operations are converted to gates using the same coarse formula as the existing analysis.
\end{enumerate}
These solver assumptions are neither introduced nor proved by this work.  They
are the methodology under which Version~2.3 claims its security levels.  Our
conclusion is therefore comparative and conditional: once the homogeneous
feature dimension is corrected and the stated rank case is used, the
candidate's own framework places the odd-characteristic sets below their
categories.

The Cabarcas--Li--Verbel--Villanueva-Polanco solver exploits additional multigraded structure in \SNOVA{} systems~\cite{Cabarcas2025SNOVA}.  We do not include any further speedup from that method.

The scaled F4 anomaly reported in \Cref{app:multibase} concerns the multibase
graph-extension presentation $B_{\rm sym}(x_1,\ldots,x_c)=w+Ky$, which
introduces linearly mixed kernel variables.  It invalidates an unqualified
production-scale extrapolation for that stronger proposed solver.  It is not an
experiment on the directly eliminated one-base residual systems costed in
\Cref{tab:main,tab:both}, and therefore neither validates nor contradicts their
semi-regular model.  The two polynomial presentations are scheme-theoretically
equivalent as solution varieties, but Gr\"obner behavior need not be invariant
under such an extended graph formulation.

\section{Implementation and Algebraic Validation}\label{sec:validation}

The statements in this section are author-reported exact computations.  They
become independently checkable only when the public ePrint is accompanied by an
immutable artifact containing the identified keys, matrices or rank
certificates, scripts, raw logs, software versions, commands, seeds, resource
limits, and a complete hash manifest.  The present \texttt{.tex} source alone
does not supply that evidence.  Subject to that artifact qualification, we
performed seven complementary layers of validation.

\paragraph{Universal identity.}
A self-contained checker verifies \Cref{lem:powerpair} over randomly sampled symmetric matrices and constructs the duplication map explicitly.  For $(m_1,\ell)=(5,4)$ the ordered feature dimension is 80, the unordered feature dimension is 50, and the duplication map has rank 50.

\paragraph{All-parameter quotient-rank certificates.}
For the simplest public relation ($R_0=I$ and $R_j=0$ for $j\ge1$), the quotient matrices attain full
column rank for all nine $q=19$ Version~2.3 shapes:
\[
\begin{array}{c|c|c}
(v,o,\ell,r)&\text{matrix size}&\text{rank}\\\hline
(28,5,4,4)&80\times50&50\\
(48,16,2,2)&64\times48&48\\
(28,4,4,5)&80\times50&50\\
(40,7,4,4)&112\times70&70\\
(72,24,2,2)&96\times72&72\\
(38,5,4,5)&100\times70&70\\
(50,9,4,4)&144\times90&90\\
(96,32,2,2)&128\times96&96\\
(52,6,4,6)&144\times90&90.
\end{array}
\]
For the Level-I $\ell=4$ constants, all $19^3=6{,}859$
scalar-coefficient column relations $(R_0,R_1,R_2,R_3)=(I,c_1I,c_2I,c_3I)$ in
the affine chart $R_0=I$ (the standard normalization; this is one chart of
$\mathbb P^3(\F_{19})$, not all $7{,}240$ projective points) produced an
$80\times50$ matrix of rank 50, and 10,000 sampled module-valued relations did
the same.  These certificates show that the universal quotient ceiling is attainable on
every Version~2.3 shape, rather than extrapolated from Level~I.  They do not by
themselves certify the quotient rank of the rejection-bypassing relation used
for a selected-target attack.  A quotient-rank defect would reduce the equation
count, but its complete effect still depends on the associated affine rank as
described in \Cref{prop:affine-rank}.

\paragraph{Official-key reconstruction.}
An independent Python implementation of the relevant Version~2.3 key-generation routines reproduced the Level-I $(28,5,19,4)$ KAT public key byte-for-byte.  This checks the field arithmetic, the fixed $A,B,Q$ generation (including the documented coefficient repair), and the public-matrix expansion used in the attack harness.

\paragraph{End-to-end verifier equivalence.}
On the KAT-derived $(28,5,4,4)$ key and two further generated keys of the same
shape, we constructed the affine-column system for one fixed public relation
with the skew offsets of \Cref{lem:skew}.  These are the instances in which the
manuscript directly certifies both the quotient and affine ranks for a
rejection-bypassing attack family.  Every trial gave
\[
 \rank(\overline E_R)=50,
 \qquad \dim\ker(\overline E_R^\transpose)=30,
 \qquad \rank(W_RL_R)=30.
\]
After eliminating the 30 affine equations, a known solution satisfied the residual 50-equation system, and reconstructing the original 80 verification coordinates gave zero discrepancy.  The same harness verified that the offset construction of \Cref{lem:skew} made symmetric-block rejection impossible for every residual assignment.

\paragraph{Cross-column normal form and restricted structure.}
The complete Level-I symmetry-reduced feature map has size $80\times680$ and rank 80.  Every individual cross-column submap is $80\times80$ and invertible, yielding \Cref{prop:crosspair}.  For 100 sampled choices of the fixed column $x_0$, the coefficient map $C_{x_0}$ had rank 80 and a 52-dimensional kernel.  Restriction of the 50 self-quadrics to those kernels gave common-radical dimension zero in every trial, individual ranks 51 or 52, and a full 50-dimensional quadratic span.  A separate adjoint-algebra computation on five keys returned only the scalar algebra.  These experiments were added specifically to test whether the 50-in-52 normal form hides an easier decomposition; none was observed.

\paragraph{Low-degree sanity check.}
For the exported official 50-in-102 residual, the degree-one Macaulay
multiplication matrix had full column rank (5,150 columns) and no unexpected
degree-one syzygy.  More substantially, the planted residual root has full
Jacobian rank 50.  Among 10,000 random coordinate 50-dimensional strips through
that root, 9,466 had invertible restricted Jacobian, closely matching the
random-matrix probability $\prod_{j=0}^{49}(1-19^{j-50})=0.9445987429\ldots$.
Across 20 independently generated official-shaped keys, every planted root again
had Jacobian rank 50, and the number of the 10,000 sampled coordinate strips
retaining full rank ranged from 9,412 to 9,491 (mean 9,447).

All 102 coordinate directions, the planted-root direction, and 10,000 random
nonzero directions produced polar matrices of full output rank 50.  We also
evaluated every projective Hessian combination whose coefficient vector has
Hamming weight at most two, $50+\binom{50}{2}(19-1)=22{,}100$ combinations,
together with 20,000 uniformly random nonzero combinations.  The minimum
observed Hessian rank was 100 in 102 variables.  These tests provide strong
negative evidence against a common radical, a low-rank pencil, or an obvious
decomposition.

They do not by themselves prove that the production systems have the semi-regular
operating degrees in \Cref{tab:actual-cores}.  We therefore ran the direct solver
audit described above: from the official KAT keys we specialized each residual to
its estimator-selected $\ell=4$ core and compared its degree-by-degree Faug\`ere
F4 behaviour (per-degree Macaulay matrix size, new basis elements, and zero
reductions) against an ensemble of five random systems matched in field,
dimensions, quadratic support density, and consistency status.  Because a
complete solve at the modeled degree is itself of the order of the estimated
attack, no run completes; the reproducible quantity is the per-degree structure
up to the maximum degree reached under a fixed memory/time budget.

For the load-bearing Level-I $\ell=4$ core ($50/41$, modeled degree 14), the
official system showed no qualitative deviation from the five matched random
controls through the maximum degree reached (F4 degree~4 completed; degree~5
entered before the budget cut both official and controls): the degree schedule
matched, no early degree fall appeared, no zero-reduction surplus beyond the F4
Koszul artifacts common to the controls was observed, and no Macaulay-rank
deficit appeared.  The official core lay within the control ensemble's spread
for the recorded quantities.  This small, truncated comparison is consistent
with semi-regular behaviour, but it is not a statistical indistinguishability
test and does not approach the modeled degree~14 solve.  The row therefore
remains a narrow estimated shortfall.  The Level-III and Level-V $\ell=4$ cores
($70/55$ and $90/73$) are larger still and lie even further beyond the
completion budget than the $50/41$ core, so no F4-audited degree is available
for them and their rows stay modeled predictions.

\paragraph{Scaling ladder.}
Because the load-bearing cores are too large to solve to completion, we validate
the operating-degree model on a ladder of \emph{faithful} symmetry-reduced cores
small enough to run msolve to a full rational parametrization.  Each core is built
through the same verified pipeline (residual construction, reduction, and
specialization), passes the identity and planted-root checks, and is kept over
$q=19$ so that field equations (active only at degree~19) never contaminate the
measured degree.  For each core we record the maximum F4 round degree and compare
against three random controls matched in dimensions, quadratic-support density,
and planted-root consistency (\Cref{tab:ladder}).

\begin{table}[t]
\centering
\caption{Observed maximum F4 round degree (msolve, run to completion) versus the
semi-regular prediction on faithful reduced cores.  $K$ is the quotient-capped
feature-map rank; each core is determined ($m=n=K$).  Daggered entries sit one
below the semi-regular first-fall index: on both, the F4 bulk table is identical
to a matching rung that reaches the predicted degree (e.g.\ the two $12/12$
cores share an identical degree schedule and peak Macaulay size), and the $\pm1$
is a trailing spair-reduction label in msolve, not a degree fall.  No audited
core reaches a degree above its prediction.  The three controls per rung are a
small qualitative comparison, not a formal statistical test.}
\label{tab:ladder}
\begin{tabular}{cccccc}
\toprule
$\ell$ & $K$ ($m/n$) & density & predicted $D_\mathrm{reg}$ & observed F4 degree & random controls \\
\midrule
2 & 3 ($3/3$)    & 0.72 & 4  & 4              & 3, 4, 3   \\
2 & 6 ($6/6$)    & 0.94 & 7  & 7              & 7, 7, 7   \\
2 & 9 ($9/9$)    & 0.95 & 10 & 10             & 5, 10, 10 \\
2 & 12 ($12/12$) & 0.96 & 13 & $12^{\dagger}$ & 7, 12, 12 \\
3 & 6 ($6/6$)    & 0.90 & 7  & 7              & 4, 6, 4   \\
3 & 12 ($12/12$) & 0.95 & 13 & 13             & 7, 7, 13  \\
4 & 10 ($10/10$) & 0.96 & 11 & $10^{\dagger}$ & 6, 10, 11 \\
\bottomrule
\end{tabular}
\end{table}

On five of the seven rungs the observed solving degree equals the semi-regular
prediction exactly; on the two daggered rungs it is one lower for the
bookkeeping reason noted in the caption.  In this ladder no audited core reaches
a degree above its prediction.  Relative to the three matched controls per
rung, the official core lies within the observed control range on six rungs and
one degree above it on one ($\ell=3$, $K=6$: $7$ versus a maximum of $6$).
This small sample is qualitatively consistent with the operating-degree model,
but it is too small for a claim of statistical indistinguishability or a
reliable estimate of run-to-run variation.  It gives direct F4 corroboration on
genuinely reduced \SNOVA{} cores only up to $K=12$; extrapolation to the
headline dimensions continues to rest on the semi-regular predictor.

The three $\ell=2$ rows of \Cref{tab:actual-cores} are the dominant determined
\emph{blocks} of a Hashimoto change-of-basis reduction; producing those exact
blocks as explicit polynomial systems is outside the current reproducibility
tooling, so their operating degrees remain the modeled semi-regular predictions
(agreed by two independent predictors) and were \emph{not} F4-audited.  As a
structural cross-check on the same official $\ell=2$ keys we instead audited the
valid but not gate-optimal square-specialization cores ($48/36$, $72/60$,
$96/79$), which are coordinate-emittable and verified byte-for-byte against the
reference public key.  They should be included in the immutable reproducibility
artifact described in \Cref{app:repro}.  These cores are themselves above the
completion budget ($n\ge36$), so they are not F4-audited; the direct degree
corroboration comes from the smaller faithful cores of \Cref{tab:ladder}.

\section{Countermeasures and Standardization Implications}\label{sec:countermeasures}

The immediate requirement is to evaluate every symmetric-matrix parameter set using
\[
 K_d=m_1\binom{d+1}{2},
 \qquad d=\dim_{\F_q}\F_q[S],
\]
as the maximum homogeneous rank in the affine-column attack; the current
irreducible construction has $d=\ell$.  Merely choosing $A,B,Q$ matrices for which the ordered-coordinate $E_R$ has high rank cannot overcome the quotient.

Possible design responses include:
\begin{enumerate}
\item \textbf{Reparameterization.} Increase the dimensions until the symmetry-reduced residual MQ system meets the desired category under the team's accepted cost methodology.  Under the rank and solver assumptions used in the tables, all three NIST categories fall short, and the Level-III and Level-V shortfalls of up to 44 bits demand larger parameters than the narrow Level-I margin alone would suggest, working against the compact keys that motivated the odd-characteristic design.
\item \textbf{Remove public-matrix symmetry.} This sacrifices the key-size reduction that motivated the odd-characteristic change, but removes \Cref{lem:powerpair} in its present form.
\item \textbf{Redesign the feature layer.} Preserving ordered power-pair coordinates requires changing the public forms or the left/right power actions so that swapping the two labels no longer gives the same quadratic polynomial.  Any such construction should then prove rank on its actual feature space.  Merely adding or randomizing more $A,B,Q$ emulsifier terms inside the present symmetric construction cannot restore the discarded alternating directions.
\item \textbf{Publish joint rank certificates.} For every proposed parameter set, evaluate both the quotient rank $\rho$ and the affine rank $\lambda_R$ on the concrete public construction and on the relations used by the attack.  A quotient-rank certificate alone does not establish selected-target solvability; an affine-rank defect must be charged through the corresponding target-consistency or resampling factor.
\end{enumerate}

The symmetric-block rejection should not be counted as a countermeasure against the affine-column attack, because fixed offsets can satisfy it identically.  A rejection rule capable of blocking the attack would need to constrain the affine family itself without destroying correct signing, rather than test a property that the attacker can set with constants.

NIST advanced \SNOVA{} to Round~3 while allowing updated specifications and implementations~\cite{NISTRound3}.  The Version~2.3 draft and Version~2.4 analysis preview are therefore at an appropriate stage for correcting the attack model.  Returning unchanged to the submitted $q=16$ construction is not an obvious remedy, because the structure-aware wedge attack already reports those parameters broken.  Any viable repair must address both attack families and be analyzed on the actual, symmetry-reduced feature space.

\section{Conclusion}\label{sec:conclusion}

The public odd-characteristic \SNOVA{} redesign makes its scalar-expanded base matrices symmetric.  That key-size optimization identifies the ordered power-pair coordinates used by the Beullens affine-column attack.  The homogeneous feature map factors through
\[
 \bigoplus_{i=1}^{m_1}\Sym^2_{\F_q}(\F_q[S])
\]
and has rank at most $m_1\binom{\ell+1}{2}$, not $m_1\ell^2$.  This is a universal algebraic identity, not a weak-key phenomenon.

Under the candidate's own semi-regular Hashimoto/Wiedemann framework, and in
the full-rank case $\rho=K$, $\lambda_R=M-K$, every $q=19$ Version~2.3 set and
every Version~2.4 preview set is estimated below its intended category.  For a simple relation, full quotient rank is certified for all nine
Version~2.3 shapes.  Simultaneous full quotient and affine rank for the
rejection-bypassing relation is certified end-to-end on the audited
$(28,5,4,4)$ Level-I shape and remains an explicit genericity assumption for
the other shapes.  The Level-III and Level-V
shortfalls are robust within a substantially wider range of solving-degree
error.  The Level-I $\ell=2$ row retains a 12.70-bit estimated margin.  The two
Level-I $\ell=4$ rows are narrower: the strict estimator costs actual $50/41$
cores at modeled degree 14, and a rise to degree 16 could erase their 4.06-bit
shortfall.  We therefore describe those rows specifically as narrow conditional
estimates rather than placing all nine rows on an equal empirical footing.

The verifier's symmetric-block rejection is bypassed constructively by fixed
skew offsets.  Independently, when a sampled fixed column gives a surjective
cross map, choosing a target-matching kernel coset restricts the official
Level-I verifier fiber to a sound 50-equation system on a 52-dimensional affine
slice.  The official Level-I residual exhibits full planted-root Jacobian rank,
full sampled polar rank, and very high sampled Hessian rank.  These facts rule
out several obvious structural degeneracies without proving the semi-regular
operating degree.

The submitted $q=16$ package is outside this exact quotient because its public
matrices are not symmetrized, but it is already heavily affected by wedge attacks.
The present result therefore targets the public odd-characteristic repair
direction rather than offering a safe fallback.  What remains uncertain is
not only solver cost: outside the audited Level-I instances the simultaneous
affine-rank and target-consistency conditions are unproved; the clean
cross-column surjectivity certificate has not been shown to coexist with the
fixed offsets used for the rejection bypass; the load-bearing $\ell=2$
Hashimoto blocks were not emitted and F4-audited as explicit systems; and the
computational certificates are not independently reproducible until an
immutable artifact is archived.  The symmetry quotient and rank ceiling are
universal.  The selected-target reduction to the displayed dimensions is exact
whenever the measured equalities $\rho=K$ and $\lambda_R=M-K$ hold.  The
cross-column normal form is an exact, sound restriction conditional on its public
surjectivity test, but not a complete parametrization of the verifier fiber.  The
displayed gate exponents additionally inherit the heuristic semi-regular
methodology used by the \SNOVA{} submission.

\appendix
\setcounter{section}{0}
\renewcommand{\thesection}{\Alph{section}}
\renewcommand{\theHsection}{app.\Alph{section}}

\section{Parameter Formulas}\label{app:parameters}

For convenience, \Cref{tab:params} records the intermediate dimensions.  Recall
\[
 n=v+o,
 \quad m_1=\left\lceil\frac{or}{\ell}\right\rceil,
 \quad M=or\ell,
 \quad K=m_1\binom{\ell+1}{2},
 \quad N_{\rm res}=n\ell-M+K.
\]

\begin{table}[ht]
\centering
\caption{Full-rank residual dimension calculation for all $q=19$ Version~2.3 draft sets ($\rho=K$, $\lambda_R=M-K$).}\label{tab:params}
\small
\begin{tabular}{@{}c l r r r r@{}}
\toprule
Level & $(v,o,\ell,r)$ & $m_1$ & $M$ & $K$ & $N_{\rm res}$\\
\midrule
I   & $(28,5,4,4)$  & 5  & 80  & 50 & 102\\
I   & $(48,16,2,2)$ & 16 & 64  & 48 & 112\\
I   & $(28,4,4,5)$  & 5  & 80  & 50 & 98\\
III & $(40,7,4,4)$  & 7  & 112 & 70 & 146\\
III & $(72,24,2,2)$ & 24 & 96  & 72 & 168\\
III & $(38,5,4,5)$  & 7  & 100 & 70 & 142\\
V   & $(50,9,4,4)$  & 9  & 144 & 90 & 182\\
V   & $(96,32,2,2)$ & 32 & 128 & 96 & 224\\
V   & $(52,6,4,6)$  & 9  & 144 & 90 & 178\\
\bottomrule
\end{tabular}
\end{table}

\section{A Linear-Algebra View of the Quotient}\label{app:linear}

For one base form, order the coordinates of $\F_q^{\ell^2}$ by pairs $(a,b)$.  Order the coordinates of $\Sym^2(\F_q^\ell)$ by pairs $a\le b$.  The duplication matrix $D_\ell$ has one $1$ in row $(a,a)$ and, for $a<b$, one $1$ in each of rows $(a,b)$ and $(b,a)$.  It has full column rank $\binom{\ell+1}{2}$ over every field.  For $m_1$ base forms,
\[
 D=I_{m_1}\otimes D_\ell.
\]
The attack matrix used in the odd-characteristic homogeneous system is $E_RD$.  In particular, even an ordered-coordinate matrix $E_R$ of full row and column rank cannot produce rank above $\rank(D)=K$.

For $\ell=4$ the single-form duplication matrix maps ten coordinates to sixteen positions.  Its six-dimensional complementary alternating subspace is spanned by the differences $e_{a,b}-e_{b,a}$ for $a<b$ (using the odd-characteristic direct-sum identification).  These six directions are always killed; additional kernel may occur after composition with the public output map.  Across $m_1=o$ square instances, the universal quotient removes $6o$ ordered coordinates.  For the Level-I $o=5$ set, the audited relation has full quotient rank and hence a 30-dimensional left kernel, exactly matching the affine-linear block observed end-to-end.

\section{Why the Extension Field Does Not Give an Immediate MinRank Break}\label{app:noAshortcut}

The field notation $A=\F_q[S]\cong\F_{q^\ell}$ makes a tempting but false
inference: one might try to regard the public matrices as an $A$-linear matrix
code of size $n\times n$ and search for an $A$-rank-one dual element.  The
structural proof above does not justify that step.  For fixed $u$, write
$\Phi_{u,i}(\xi,\eta)=\Psi_i(\xi\otimes\eta)(u)$.  This pairing is
symmetric and bilinear over $\F_q$, but it is generally not $A$-balanced:
\[
 \Phi_{u,i}(\alpha\xi,\eta)
 \ne
 \Phi_{u,i}(\xi,\alpha\eta).
\]
Such balance would require the public form to commute appropriately with the
$A$-action.  In \SNOVA{}, the secret transform and the $Q$ factors lie in
$A$, but the base quadratic blocks are sampled as general matrices subject
only to the UOV and, in odd characteristic, global symmetry conditions.  The
specification itself treats the extension-field decomposition as a collection
of Frobenius-conjugate lifted UOV instances, not as one $A$-linear public
matrix code.

An official-KAT audit of the Level-I $(28,5,19,4)$ key gives a particularly
sharp falsification.  With $J=I_n\otimes S$, all five public matrices satisfy
\[
 \rank(P_iJ-JP_i)=132=n\ell.
\]
None of their $5\cdot33^2=5{,}445$ scalar $4\times4$ blocks lies in the
centralizer field $A$.  The two-sided power orbit has the full tensor-square
dimension $5\ell^2=80$, while its symmetric quotient has the full dimension
$5\binom{\ell+1}{2}=50$; it does not collapse further to the
$5\ell=20$ dimensions that an $A$-balanced model would predict.  Finally, a
genuine public oil dual $uu^\transpose$ has $\F_q$-rank one, fails to commute
with $J$, and has nonzero blocks of rank one.  A nonzero multiplication matrix
from the field $A$ is invertible, so this dual is not an $A$-valued matrix.

Hence there is no justified $n\times n$ rank-one MinRank problem over
$\F_{q^\ell}$ that replaces the scalar-expanded problem.  The valid surviving
uses of the extension field are the known Frobenius/diagonal-block lifting
attacks and the $A$-valued Gram relations analyzed above.

\section{The Multibase Stable-Ideal Extension}\label{app:multibase}

\Cref{prop:cbase} gives an exact attack architecture beyond the one-base
reduction.  Specializing to the current irreducible degree-$\ell$ construction,
for a $c$-base relation let
\[
 F_c=m_1\binom{c\ell+1}{2},\qquad M=or\ell,
 \qquad p_c=F_c-M.
\]
When $F_c\ge M$ and the public feature map is surjective, select a preimage
$w$ of the target and a basis $K$ of its kernel.  Forgery becomes
\begin{equation}\label{eq:lifted-multibase}
 B_{\rm sym}(x_1,\ldots,x_c)=w+Ky,
\end{equation}
with $p_c$ linear kernel variables.  For $c=r$, this is a symmetry-aware
stable encoding of direct forgery.  Surjectivity makes every target lie in a
linear output fibre, so no separate \emph{linear} membership search is needed;
it does not imply that the nonlinear Gram map has a preimage in that fibre.

After diagonalizing the public $A$-action, there remain $\ell$ grading blocks,
not $c\ell$ blocks.  If each block contains $s=M/\ell=or$ specialized
quadratic variables, the numbers of feature equations of each multidegree are
\[
 A_c=m_1\binom{c+1}{2}
 \quad\text{of degree }2e_a,
 \qquad
 B_c=m_1c^2
 \quad\text{of degree }e_a+e_b.
\]
The corresponding formal complete-intersection series is
\[
 H(\mathbf t)=
 \frac{\prod_a(1-t_a^2)^{A_c}
       \prod_{a<b}(1-t_at_b)^{B_c}}
      {\prod_a(1-t_a)^s}.
\]
Using the Cabarcas-style nonhomogeneous criterion gives attractive formal
screens.  For the Level-I square $\ell=4$ set, $c=4$ gives about
$2^{99.8}$ operations with an $\omega=2$ sparse-linear-algebra model and
$2^{132.3}$ with $\omega=2.81$.  These figures are \emph{not} part of the
headline attack claim.  They depend on a solving-degree extrapolation that has
not been validated at production scale.

Scaled experiments reveal the risk.  Generic lifted systems with
$(\ell,c,m_1)=(4,4,1)$ reached maximum F4 degrees 4 and 6 when the number of
variables per grading block was one and two, respectively; the next scale
already produced very large degree-three matrices before completion.  A
separate $(3,3,1)$ family reached degree 6 where the formal criterion predicted
roughly 3--4.  These are ordinary-F4 experiments rather than a custom
multigraded implementation, so they neither prove nor refute the attack
estimate.  They do show that a public complexity claim near $2^{100}$ would be
premature without an independent multigraded-F4 or Magma validation.

The graph extension is scheme-theoretically equivalent to direct forgery:
projection onto the $x$-coordinates identifies its solution variety with the
corresponding fiber of the direct feature map.  It should not, however, be
confused with the directly eliminated one-base polynomial presentation.  The
public kernel generally mixes the nominal grading sectors, so the graph-extension
variables do not decompose into the independent homogeneous blocks assumed by the
most optimistic formal complete-intersection screen.  Gr\"obner behavior can
therefore differ substantially even though the two presentations encode the same
solutions.

The exact conclusions are therefore: the multibase factorization, feature counts,
graph equivalence, and surjectivity tests are valid; the scaled F4 anomaly applies
to this extended formulation; and the stronger multibase work factors remain an
unconfirmed solver extrapolation.  It supplies no direct evidence for or against
the one-base estimator cores in \Cref{tab:actual-cores}.

\section{Reproducibility and Claim Checklist}\label{app:repro}

A public claim based on this work should distinguish the following statements.

\begin{description}[leftmargin=0pt]
\item[Proved algebraically.] The homogeneous affine-column system factors through unordered power pairs and has rank at most $m_1\binom{d+1}{2}$, with $d=\dim\F_q[S]$; the verifier's symmetric-block rejection can be bypassed by explicitly computable fixed skew offsets; and \Cref{thm:residual} is an exact solution-preserving reduction conditional on target consistency.
\item[Author-reported computations.] The manuscript reports that the Level-I official-key reconstruction matches a KAT public key; one rejection-bypassing relation on the KAT-derived $(28,5,4,4)$ key and two generated keys of that shape has $\rho=50$ and $\lambda_R=30$ and yields an exact 50-in-102 system; all six official Level-I cross-feature output maps are invertible, and 100 sampled fixed columns pass the surjectivity test needed for the sound 50-in-52 slice; a separate simple relation has full quotient-rank certificates on all nine Version~2.3 shapes; and the official Level-I residual has full sampled polar rank, full planted-root Jacobian rank, and high-rank sampled Hessian combinations.  These become independently certified only through the immutable artifact specified below.
\item[Assumed for untested shapes.] The all-shape selected-target dimensions require the same attack relation and offsets to satisfy $\rho=K$ and $\lambda_R=M-K$.  These simultaneous equalities are directly certified only on the audited $(28,5,4,4)$ Level-I shape; a defect elsewhere would alter the residual dimensions or require the target-consistency factor in \Cref{prop:affine-rank}.
\item[Estimated heuristically.] The gate complexities in \Cref{tab:main,tab:both}, like the submission's estimates, assume that the sampled residual fiber is solvable with the modeled probability, semi-regular behavior, the modeled success of guessing and specialization, and a particular sparse-linear-algebra model.  The two Level-I $\ell=4$ rows are sensitive to a one- or two-degree error; the Level-III and Level-V margins tolerate several additional degrees in the same model.  The load-bearing $\ell=2$ Hashimoto blocks are prediction-only and were not explicitly generated or F4-audited.
\item[Audited negative result.] The public forms are not $A$-linear and there is no immediate $n\times n$ MinRank shortcut over $\F_{q^\ell}$.
\item[Preliminary extension.] The exact multibase Gram attack is included as a solver-research direction, but its stronger formal exponents are not claimed as validated production work factors.
\item[Not claimed.] We do not claim a practical production-size forgery, a key-recovery attack, a break of the $q=16$ parameter sets, a universal full-affine-rank theorem, a proof that every restricted target fiber is nonempty, or a proven running-time upper bound matching the estimates.
\end{description}

The public ePrint should archive an immutable reproducibility artifact at the
same time as the manuscript.  At minimum, that archive must contain:
\begin{itemize}
\item exact repository commits, local patches, parameter files, KAT identifiers, and generated public-key digests;
\item a self-contained checker for the symmetry identity, quotient map, parameter formulas, and rejection bypass;
\item the concrete matrices or machine-verifiable rank certificates for every quotient-rank, affine-rank, and cross-map claim, together with a verifier;
\item a pinned implementation of both semi-regular conventions, the Hashimoto optimization, the multibase Hilbert-series calculation, and all exact cost tables;
\item rank-search, Jacobian, coordinate-strip, polar, Hessian, and degree-one Macaulay scripts and raw logs;
\item the solving-degree experiments, including emitted polynomial systems, matched controls, solver versions, commands, seeds, resource limits, and complete per-degree F4 logs;
\item end-to-end KAT and verifier-equivalence transcripts; and
\item a manifest giving the SHA-256 digest of every source, input, output, and evidence file, plus a documented one-command reproduction path.
\end{itemize}
Until such an archive is linked, the computational statements above should be
read as author-reported rather than independently reproduced.  A hash of the
specification PDF alone is useful but insufficient to identify the code, keys,
certificates, logs, and local modifications.  The SHA-256 digest of the
Version~2.3 PDF used in the implementation audit is
\begin{center}
\ttfamily\small 963fa43d7e6ed8ecf1e77556e7cf47f62d4c34177efece2cd9a9517d0fbaf783.
\end{center}

\printbibliography

\end{document}
